% Compile twice: lualatex 30003989.tex
\documentclass[11pt]{article}
\usepackage[margin=25mm]{geometry}
\usepackage{fontspec,unicode-math}
\setmainfont{Latin Modern Roman}
\setmathfont{Latin Modern Math}
\usepackage{amsmath,amssymb,mathtools}
\usepackage{xurl,hyperref}
\hypersetup{colorlinks=true,urlcolor=blue,linkcolor=blue}
\setcounter{secnumdepth}{0}
\setlength{\parindent}{0pt}
\setlength{\parskip}{5pt}
\setlength{\emergencystretch}{3em}
\newcommand{\sref}[2]{\hyperlink{source:#1:#2}{[S#2]}}
\newcommand{\cataloguescope}[1]{\paragraph{Source question.}#1}
\newcommand{\floor}[1]{\left\lfloor#1\right\rfloor}
\begin{document}
% UnsolvedMath ID 30003989; source catalogue code OWR-16633-003
\setcounter{section}{30003988}
\section{Proximity and Sparsity of Optimal Integer Solutions}\label{30003989}
{\small\sffamily UnsolvedMath ID 30003989 · Computer Science}
\subsection{Definitions and mathematical statement}

\paragraph{Shared notation.}Standard number systems are denoted by $\mathbb N,\mathbb Z,\mathbb Q,\mathbb R,\mathbb C$; index ranges, computational models and specialized notation follow the cited problem source and the definitions below.


For a vector $z$, $\lVert z\rVert_1=\sum_i|z_i|$, $\lVert z\rVert_\infty=\max_i|z_i|$, and sparsity is the number of nonzero coordinates. Proximity measures the distance between a continuous optimum and an integral optimum.

A relaxation contains the integral feasible solutions. For a nonnegative minimization objective, its integrality gap is the supremum of the ratio between the integral optimum and the relaxation optimum over the specified instances.

A linear program optimizes $c^Tx$ over a polyhedron $\{x:Ax\leq b\}$. In the simplex method a pivot moves between adjacent basic feasible solutions; a pivot rule chooses the next such move.

An integer program imposes integral coordinates on a system of constraints. A mixed-integer program allows both integral and real coordinates; $\operatorname{conv}(S)$ denotes the set of finite convex combinations of feasible points of $S$.
\paragraph{Statement recorded in the supplied source.}
For integer linear optimization, given an optimal solution of the linear-programming relaxation, how close can an optimal integer solution be in the $\ell_1$ or $\ell_\infty$ norm, and for a system of linear equations with nonnegative variables and a linear objective, how sparse can an optimal integer solution be? \sref{30003989}{2}
\paragraph{Scope qualification.}
The matrix parameters allowed in the bounds and the optimal-solution selection quantifiers are source-dependent.
\paragraph{Source status.}
The supplied catalogue labels this record partially solved. Its source-status metadata was checked against the supplied snapshot on 2026-10-11. This import does not independently establish current open status or mathematical correctness. The catalogue literature assessment is dated 2026-08-21.
\subsection{Short English statement}
Bound distance from linear-programming optima to integer optima and the number of nonzero coordinates needed by integer optima.
\subsection{Sources}
\begin{itemize}
\item[\hypertarget{source:30003989:1}{S1}] Ulam AI and the UnsolvedMath Contributors, UnsolvedMath: A Curated Collection of Open Mathematics Problems, record 30003989 (OWR-16633-003). Catalogue curation and metadata: CC BY 4.0. The supplied source question is retained; definitions and the English summary are editorial additions.
\url{https://huggingface.co/datasets/ulamai/UnsolvedMath}.
\item[\hypertarget{source:30003989:2}{S2}] Combinatorial Optimization (2019), Oberwolfach Reports, DOI 10.4171/owr/2018/50, pp. 2969--3023
\url{https://doi.org/10.4171/owr/2018/50}.
\item[\hypertarget{source:30003989:3}{S3}] The Support of Integer Optimal Solutions
\url{https://doi.org/10.1137/17M1162792}.
\item[\hypertarget{source:30003989:4}{S4}] Optimal Proximity Bound and Product Function Estimates in Integer Linear Programming
\url{https://arxiv.org/abs/2606.13579}.
\textit{Underlying Oberwolfach Reports material retains its source-specific licence, including CC BY-SA 4.0 where applicable; see THIRD\_PARTY\_NOTICES.md.}
\end{itemize}
\paragraph{Preparation and review.}
Original question attributed to the cited source and UnsolvedMath contributors. Codex (GPT-6; exact serving revision not exposed) assisted with definitions, English summaries, source inspection and formatting. No new human mathematical authorship or independent proof verification is asserted. Canonical ID 30003989 and source code OWR-16633-003 checked against the complete supplied registry and existing repository statements; known flagged repetitions and obvious identical conjectures excluded. Full mathematical-equivalence/literature review has not been performed. Statement preparation and final review followed contribution prompts C and D; automated checks verify structure and identity, not open status.
\label{end:30003989}
\end{document}
